\section{Conclusion}
\label{sec:conclusion}

We have introduced \ilp{} redistricting, the first structure-layer algorithm
in this platform to provide a mathematical optimality certificate for
edge-cut compactness.
For small instances ($n \leq 500$ tracts), the ILP formulation with
both-endpoint MTZ contiguity constraints and a single cut-edge indicator
per undirected edge finds the provably optimal redistricting plan in
practical solve times: under 50\,s for NH ($n=260$, $k=2$) and under
1\,s for single-district states.

\paragraph{What ILP is the right tool for.}
\ilp{} redistricting is the correct choice in two situations:

\begin{enumerate}
\item \textbf{Small-state certification.}
  For VT, ND, WY, AK, DE, SD, MT ($k=1$) and NH, RI ($k=2$), \ilp{}
  is strictly superior to any heuristic: it finds a better plan (lower EC)
  in competitive or better time, and adds the certificate that no heuristic
  can provide.

\item \textbf{Bisection node certification.}
  For large-$k$ states where each bisection node has $n \leq 500$ tracts,
  running \ilp{} at every node certifies the compactness of each sub-partition
  independently of tree structure.
  For $k = 14$ (NC), depth-3 nodes have $n \approx 330$ tracts; \ilp{}
  is tractable at these nodes.
\end{enumerate}

\paragraph{What ILP is not the right tool for.}
For production 50-state builds, or any state where bisection nodes exceed
500 tracts, \metis{} remains the correct default.
The ILP size guard (default \texttt{max\_tracts = 500}) falls back to
\metis{} automatically and records the fallback in the audit chain.

\paragraph{Phase~2 enhancements.}
Several enhancements are deferred:

\begin{enumerate}
\item \textbf{DFJ branch-and-cut.}
  Replacing MTZ with Dantzig--Fulkerson--Johnson cuts~\citep{nemhauser1988}
  yields a tighter LP relaxation and typically faster branch-and-bound.
  Requires a lazy-constraint callback in the solver interface.
  Deferred to Phase~2.

\item \textbf{\highs{} full integration.}
  Phase~1 uses \highs{} as a subprocess.
  Phase~2 will use the \texttt{highs-sys} Rust bindings for tighter
  integration (in-process solving, MIP start hints for warm-start from
  \metis{}).

\item \textbf{Gurobi and CPLEX backends.}
  The \texttt{Solver} trait allows new backends without modifying the
  formulation.
  Commercial solvers (Gurobi, CPLEX) are $2$--$5\times$ faster than
  \glpk{} on MIP instances and would extend the tractable range to
  $n \approx 1{,}000$.

\item \textbf{Warm-start from \metis{}.}
  Providing the \metis{} solution as an MIP start hint reduces the
  gap between the initial upper bound and the LP lower bound,
  often halving solve time.
  Requires checking \glpk{}/\highs{} support for MIP start APIs.

\item \textbf{Partial certification for large states.}
  For states where some bisection nodes exceed 500 tracts, run \ilp{}
  at the tractable nodes and \metis{} at the larger nodes.
  Record which nodes are certified and which are heuristic in the audit
  chain.
  This gives a ``partially certified'' plan with precise documentation
  of where the certification is and is not in force.
\end{enumerate}

\paragraph{Legal caveat.}
The ILP certificate establishes mathematical optimality for edge-cut
compactness under the specified population balance tolerance and census
data vintage.
It is not a substitute for legal counsel or a certification of compliance
with state constitutional compactness standards, Voting Rights Act
requirements, or other legal criteria applicable in a given jurisdiction.
Practitioners should consult legal counsel before using the ILP
certificate in redistricting litigation or regulatory proceedings.

\paragraph{Summary.}
\ilp{} redistricting fills the certification gap in the algorithm
portfolio.
Phase~1 (this paper) provides proven-optimal plans for small instances
and bisection nodes using \glpk{} with both-endpoint MTZ contiguity.
Phase~2 will extend coverage with DFJ cuts and faster solver backends.
Together with \metis{} (heuristic, production-scale) and
\be{}~\citep{dellaLibera2026comparison} (heuristic, high-quality search),
the platform now covers the full quality-runtime-certification spectrum:
fast approximate (\metis{}), best-effort approximate (\be{}), and
exact certified (\ilp{}).
